\documentclass[11pt]{article}
\usepackage[utf8]{inputenc}
\usepackage[margin=1in]{geometry}
\usepackage{booktabs}
\usepackage{hyperref}
\usepackage{enumitem}
\usepackage{parskip}

\title{Comprehensive Progress Report: \\ Sequential Credit Risk Modeling (Jan -- Apr 2026)}
\author{Attabra Benjamin Ekow}
\date{April 11, 2026}

\begin{document}

\maketitle

\textbf{Project:} Sequential Deep Learning for Credit Risk Modeling in Data-Constrained Environments \\
\textbf{Period:} January 1, 2026 -- April 11, 2026 \\
\textbf{Status:} Framework complete; Preliminary benchmark results finalized.

\hr

\section{Executive Summary: The Path to a Rigorous Framework}
Over the past four months, the project has evolved from a basic data generation script into a comprehensive research framework. We navigated through critical data integrity challenges and model performance bottlenecks to establish a benchmark that demonstrates the high value of behavioral feature engineering over traditional sequential modeling. Our central finding is that well-engineered static features capture the primary default signal in the Ghanaian mobile money context, significantly outperforming standard sequential architectures.

\section{Detailed Progression Timeline}

\subsection{Phase 1: Foundation \& Data Synthesis (January -- February)}
\begin{itemize}[noitemsep]
    \item \textbf{The Problem:} Existing credit risk datasets are either private or lack the temporal granularity of mobile money (MoMo).
    \item \textbf{The Progression:}
    \begin{itemize}
        \item \textbf{Modular Architecture (Jan 25):} Restructured the repository from a collection of scripts into a modular package (\texttt{src/seqcredit\_model}).
        \item \textbf{Generator Rewrite (Jan 28):} Re-engineered the \texttt{CalibratedMoMoDataGenerator} from a global generator to a per-user simulator, enabling the creation of individualized transaction histories and loan archetypes for 10,000 users.
        \item \textbf{Baseline Validation (Feb 17):} Established initial benchmarks using Logistic Regression and a basic LSTM to verify that transaction sequences could be processed end-to-end.
    \end{itemize}
\end{itemize}

\subsection{Phase 2: Debugging, Expansion \& The ``Imbalance Pivot'' (March)}
\begin{itemize}[noitemsep]
    \item \textbf{The Problem:} Early models showed poor performance (AUC $\sim$0.53), and label assignments were inconsistent.
    \item \textbf{The Progression:}
    \begin{itemize}
        \item \textbf{Label Integrity Fix (Mar 11):} Identified and resolved a critical bug where models were training on outdated summary files. Standardized on \texttt{user\_labels.csv} as the definitive ground truth.
        \item \textbf{Model Diversification (Mar 17):} Expanded the suite from 2 to 6 models, adding \textbf{Random Forest, XGBoost, LightGBM, and a Hybrid LSTM.}
        \item \textbf{Imbalance Correction (Mar 21):} Applied \texttt{class\_weight='balanced'} and \texttt{scale\_pos\_weight} to address the 11\% default rate. This led to a major performance jump, with static model AUCs rising from \textbf{$\sim$0.53 to $\sim$0.88}.
    \end{itemize}
\end{itemize}

\subsection{Phase 3: Interpretability, Calibration \& Trust (Late March -- Early April)}
\begin{itemize}[noitemsep]
    \item \textbf{The Problem:} High AUC alone is insufficient for credit risk; we needed to know \textit{why} models predict default and if their probabilities are reliable.
    \item \textbf{The Progression:}
    \begin{itemize}
        \item \textbf{XAI Integration (Mar 18):} Implemented SHAP analysis and Permutation Importance. Identified that \textbf{Behavioral Diversity} and \textbf{Loan History} are the primary drivers of risk.
        \item \textbf{Quantified Ablation (Apr 9):} Completed feature-group ablation studies. Dropping \texttt{behavioural\_diversity} caused the largest AUC drop (-0.09), confirming its role as the dominant signal.
        \item \textbf{Calibration Layer (Mar 30):} Introduced reliability diagrams. Found that while the Hybrid LSTM is discriminative, Tree Ensembles (RF/XGB) provide more reliable probability scores for operational use.
    \end{itemize}
\end{itemize}

\subsection{Phase 4: Scientific Rigor \& Benchmark Automation (April)}
\begin{itemize}[noitemsep]
    \item \textbf{The Problem:} Results needed statistical validation to support publication claims.
    \item \textbf{The Progression:}
    \begin{itemize}
        \item \textbf{Benchmark Automation (Apr 2):} Transitioned to dedicated evaluation scripts (\texttt{run\_cv\_benchmark.py}) using \textbf{5-Fold Stratified Cross-Validation}.
        \item \textbf{Statistical Significance (Apr 7):} Integrated bootstrap significance testing to calculate 95\% Confidence Intervals, ensuring our model comparisons are statistically robust.
        \item \textbf{Conclusion on Sequences:} Standard sequential models (LSTMs) consistently underperformed compared to well-engineered static features, suggesting that aggregated behavior is more predictive of default than raw transaction order in this benchmark.
    \end{itemize}
\end{itemize}

\section{Technical Evolution: Codebase Maturity}
\begin{table}[h]
\centering
\begin{tabular}{@{}lll@{}}
\toprule
\textbf{Feature} & \textbf{January State} & \textbf{April State} \\ \midrule
Data Source & Raw CSVs & Calibrated simulator (10k users, 5 archetypes). \\
Model Suite & 2 models (LR, LSTM) & 6 models (including Hybrid LSTM). \\
Validation & Single-split (train/test) & 5-Fold CV + Bootstrap Significance Tests. \\
Explainability & None & SHAP, Surrogate Trees, and Feature Group Ablation. \\
Performance & AUC $\sim$0.53 (Random) & AUC $\sim$0.88 (Tree Ensembles) / $\sim$0.85 (Hybrid LSTM). \\ \bottomrule
\end{tabular}
\end{table}

\section{Key Performance Benchmarks (Finalized)}
\begin{table}[h]
\centering
\begin{tabular}{@{}lll@{}}
\toprule
\textbf{Model} & \textbf{Target: Default ($y_{default}$)} & \textbf{Target: Late/Bad ($y_{bad}$)} \\ \midrule
\textbf{Random Forest} & \textbf{0.884 $\pm$ 0.014} & \textbf{0.727 $\pm$ 0.017} \\
Logistic Regression & 0.914 $\pm$ 0.007 & 0.725 $\pm$ 0.021 \\
Hybrid LSTM & 0.850 $\pm$ 0.033 & 0.705 $\pm$ 0.018 \\
Standalone LSTM & 0.523 $\pm$ 0.041 & 0.533 $\pm$ 0.015 \\ \bottomrule
\end{tabular}
\end{table}

\textbf{Core Conclusion:} Aggregated behavioral features (specifically diversity and history) capture the relevant predictive signal; the temporal order extracted by standard LSTMs adds negligible value for default prediction in this environment.

\section{Strategic Next Steps}
\begin{enumerate}
    \item \textbf{Manuscript Finalization (Apr 15):} Finalize the abstract for Deep Learning Indaba 2026 highlighting the ``Static vs. Sequence'' findings.
    \item \textbf{External Validation (Paper B):} Prepare the framework for ingestion of real Telecel Ghana transaction logs to verify if these synthetic findings hold in production.
    \item \textbf{Attention Mechanisms:} Future work will investigate if advanced attention-based sequential models can extract signals that LSTMs cannot.
\end{enumerate}

\end{document}
